\documentclass[titlepage,twoside,openright]{jsreport}
%\usepackage{thesisj}
\usepackage[dvipdfmx]{graphicx,xcolor}
\usepackage{subcaption}
\usepackage{amsmath,amssymb}
\usepackage{booktabs}
\usepackage{tabularx}
\usepackage{url}
\usepackage[dvipdfmx]{hyperref}
\usepackage{multirow}
\usepackage[cc]{titlepic}
\usepackage{fancyhdr}
\usepackage{makeidx}
\bibliographystyle{jplain}

\makeatletter
\renewcommand{\chapter}{%
  \if@openleft\cleardoublepage\else
  \if@openright\cleardoublepage\else\clearpage\fi\fi
  \global\@topnum\z@
  \secdef
      {\@omit@numberfalse\@chapter}%
      {\@omit@numbertrue\@schapter}}
\makeatother

\pagestyle{fancy}
\lhead[\thepage]{\rightmark}
\chead[]{}
\rhead[\leftmark]{\thepage}
\lfoot[]{}
\cfoot[]{}
\rfoot[]{}
\renewcommand{\headrulewidth}{1pt}
\renewcommand{\footrulewidth}{0pt}
%\fancyhead{}
%\renewcommand{\chaptermark}[1]{\markboth{第\ \thechapter\ 章\,\, #1}{}}
%\fancyhead[RO]{{\small \rightmark}}
%\fancyhead[LE]{{\small \leftmark}}
%\renewcommand{\headrulewidth}{0pt}
%\fancyfoot{}
%\fancyfoot[LE,RO]{{\small \thepage}}
%\cfoot{}

\makeindex
\title {修士論文}
\author{土井 隆暢}
\date{}
\titlepic{\includegraphics[width=0.5\columnwidth]{eps/E-C.eps}}


\begin{document}
\maketitle

\tableofcontents
\pagenumbering{roman}
\listoffigures
\pagenumbering{Roman}

\pagenumbering{arabic}
\chapter{Introduction}
\section{トリプルアルファ反応}


\chapter{Experimental setup}
\section{OKTAVIAN}
\begin{figure}
  \centering
  \includegraphics[width=0.5\columnwidth]{eps/E-C.eps}
  \caption{simbol}
  \label{fig:simbol}
\end{figure}
\ref{fig:simbol}
\section{検出器セットアップ}
\subsection{検出ガスの決定}
$\alpha$粒子のレンジで決めた。
efficiency が一番大きくなるように。
\subsection{ドリフトスピード}
ドリフト速度の決定方法は30 degree 方向に$\alpha$線源から$\alpha$を出して、
その飛跡がデータ上でどう見えるかで決定する。
ドリフト速度の時間依存性も見た。
\chapter{Analysis}
\subsection{機械学習}
これまではHough 変換を使って解析を行ってきたが、
高速に処理をするためにニューラルネットワークを用いた解析方法を開発した。
\subsection{解析}
\chapter{Conclusion and discussion}
\chapter{Acknowledge}

\bibliography{mybibfile}
\end{document}
